% !Mode:: "TeX:UTF-8"

\section{项目简介}

随着物联网技术和可穿戴设备市场的快速发展，智能手表作为一种集时间显示、运动监测、数据通信于一体的便携设备，已经成为现代生活中不可或缺的电子产品。传统手表功能单一，仅能提供基础的时间信息，难以满足用户对健康监测、运动记录、数据交互等多方面的需求。将嵌入式微控制器应用于可穿戴设备领域，能够有效提升产品的智能化水平和用户体验。

STM32系列微控制器以其高性能、丰富的外设资源（如GPIO、定时器、I2C、USART等）、低功耗特性以及成熟的开发生态，成为嵌入式系统开发的理想选择。本项目依托嵌入式系统课程，旨在设计并实现一个基于STM32F103C8T6微控制器的智能手表系统，将嵌入式软硬件设计技术应用于可穿戴设备场景。

本设计的主要功能包括：
\begin{enumerate}
    \item \textbf{OLED显示与界面交互：} 驱动SSD1306 128$\times$64 OLED显示屏，实现大字时钟、传感器数据、蓝牙状态、步数统计等多页面循环切换的用户界面。
    \item \textbf{六轴运动传感：} 通过MPU6050传感器采集三轴加速度与三轴陀螺仪数据，计算姿态角（Pitch/Roll）并提供运动数据源。
    \item \textbf{HC-05蓝牙通信：} 实现基于USART的蓝牙串行通信，设计自定义帧协议（含STX、CMD、DATA、XOR校验、ETX字段），支持手表与手机之间的数据双向传输。
    \item \textbf{计步器算法：} 基于加速度向量幅值（Acceleration Vector Magnitude）与动态阈值检测，配合最小步间间隔防抖，实现步数统计、距离估算与卡路里消耗计算。
    \item \textbf{Android上位机：} 开发Android APP，通过蓝牙串口协议连接手表，实时显示传感器数据、同步系统时间、记录通信日志。
    \item \textbf{系统集成与实践验证：} 完成从硬件焊接、I2C总线驱动、中断处理到多模块协同的完整系统设计与集成，在真实硬件平台上进行功能测试与性能验证。
\end{enumerate}

本项目不仅涵盖了嵌入式软件开发（C语言）、硬件接口驱动（OLED、MPU6050、HC-05）、数字信号处理（计步算法、姿态解算）等核心知识点，还涉及硬件电路设计（电源管理、引脚分配、I2C总线恢复）与系统调试等实践技能，是对嵌入式系统课程所学知识的综合应用与实践。